% Mathématiques 5e — cours V3
% Transformations, symétrie centrale et angles

\section{Objectifs}

\begin{itemize}
\item Définir le demi-tour ou symétrie centrale.
\item Construire l'image d'un point ou d'une figure par demi-tour.
\item Connaître les propriétés conservées par un demi-tour.
\item Reconnaître angles opposés par le sommet, adjacents, supplémentaires, alternes-internes et correspondants.
\item Utiliser des relations d'angles pour déterminer des mesures.
\item Caractériser le parallélisme à l'aide des angles alternes-internes ou correspondants.
\item Justifier une conclusion géométrique sans se fier uniquement au dessin.
\end{itemize}

\section{1. Du demi-tour à la symétrie centrale}

Un demi-tour correspond à une rotation de :
\[
180^\circ
\]
autour d'un point appelé centre.

\begin{definition}
Le point $A'$ est l'image du point $A$ par le demi-tour de centre $O$ lorsque $O$ est le milieu du segment $[AA']$.
\end{definition}

Autrement dit :
\[
OA=OA'
\]
et les points :
\[
A,\ O,\ A'
\]
sont alignés.

\coursefigure{/images/mathematiques/5eme/demi-tour.svg}{Deux triangles images l'un de l'autre par un demi-tour autour du centre O.}{Un demi-tour place chaque point image de l'autre côté du centre, à la même distance.}

\section{2. Construire l'image d'un point}

Pour construire $A'$, image de $A$ par le demi-tour de centre $O$ :
\begin{enumerate}
\item tracer la droite $(AO)$ ;
\item prolonger cette droite au-delà de $O$ ;
\item reporter la longueur $OA$ de l'autre côté de $O$ ;
\item placer $A'$ tel que $OA'=OA$.
\end{enumerate}

Alors :
\[
O
\]
est le milieu de :
\[
[AA'].
\]

\section{3. Construire l'image d'une figure}

Pour transformer un polygone :
\begin{enumerate}
\item construire l'image de chaque sommet ;
\item relier les sommets images dans le même ordre.
\end{enumerate}

La construction point par point permet de contrôler la transformation.

\section{4. Propriétés conservées}

Le demi-tour conserve :
\begin{itemize}
\item les longueurs ;
\item les mesures d'angles ;
\item l'alignement ;
\item le parallélisme ;
\item les aires ;
\item la forme des figures.
\end{itemize}

Si $A'$ et $B'$ sont les images de $A$ et $B$ :
\[
A'B'=AB.
\]

\section{5. Image d'une droite}

L'image d'une droite ne passant pas par le centre $O$ est une droite parallèle.

Si une droite passe par $O$, son image est elle-même.

Cette propriété permet de construire certaines images sans transformer tous les points.

\section{6. Angles : vocabulaire de base}

On entretient le vocabulaire suivant :
\begin{itemize}
\item angle nul ;
\item angle aigu ;
\item angle droit ;
\item angle obtus ;
\item angle plat ;
\item angle plein ;
\item angles adjacents ;
\item angles opposés par le sommet ;
\item angles supplémentaires.
\end{itemize}

Les mesures de référence sont :
\[
90^\circ,\qquad180^\circ,\qquad360^\circ.
\]

\section{7. Angles opposés par le sommet}

Lorsque deux droites se coupent, les angles opposés par le sommet ont la même mesure.

\begin{propriete}
Si deux angles sont opposés par le sommet, alors leurs mesures sont égales.
\end{propriete}

Par exemple, si l'un mesure :
\[
68^\circ,
\]
son opposé mesure :
\[
68^\circ.
\]

\section{8. Angles supplémentaires}

Deux angles supplémentaires ont une somme égale à :
\[
180^\circ.
\]

Si :
\[
\alpha=47^\circ,
\]
l'angle supplémentaire vaut :
\[
180^\circ-47^\circ=133^\circ.
\]

\section{9. Deux droites coupées par une sécante}

Lorsque deux droites sont coupées par une troisième droite appelée sécante, plusieurs paires d'angles apparaissent.

On distingue notamment :
\begin{itemize}
\item les angles alternes-internes ;
\item les angles correspondants.
\end{itemize}

\coursefigure{/images/mathematiques/5eme/angles-paralleles.svg}{Deux droites parallèles coupées par une sécante avec angles alternes-internes et correspondants repérés.}{Les égalités d'angles fournissent des critères pour reconnaître ou exploiter le parallélisme.}

\section{10. Angles alternes-internes}

Deux angles alternes-internes sont situés :
\begin{itemize}
\item entre les deux droites ;
\item de part et d'autre de la sécante.
\end{itemize}

Lorsque les deux droites sont parallèles, les angles alternes-internes correspondants ont la même mesure.

\begin{propriete}
Si deux droites sont parallèles, alors deux angles alternes-internes formés par une même sécante sont égaux.
\end{propriete}

\section{11. Angles correspondants}

Deux angles correspondants occupent la même position relative aux deux intersections.

Lorsque les droites sont parallèles :
\[
\text{angles correspondants}=\text{même mesure}.
\]

\section{12. Caractériser le parallélisme}

Le programme de 5e ne demande pas seulement d'utiliser le parallélisme connu.

Il demande aussi de le \textbf{caractériser} grâce aux angles.

\begin{propriete}
Si deux angles alternes-internes déterminés par deux droites et une sécante ont la même mesure, alors les deux droites sont parallèles.
\end{propriete}

De même, l'égalité d'angles correspondants permet de conclure au parallélisme.

\section{13. Exemple développé : parallèles connues}

Deux droites $(d_1)$ et $(d_2)$ sont parallèles.

Une sécante forme un angle de :
\[
52^\circ
\]
avec $(d_1)$.

L'angle alterne-interne correspondant avec $(d_2)$ mesure également :
\[
\boxed{52^\circ}.
\]

Aucune mesure au rapporteur n'est nécessaire : on utilise une propriété.

\section{14. Exemple développé : prouver le parallélisme}

Deux droites sont coupées par une sécante.

Deux angles correspondants mesurent :
\[
73^\circ
\]
chacun.

Comme ces angles correspondants sont égaux, on peut conclure :
\[
\boxed{(d_1)\parallel(d_2)}.
\]

La conclusion est obtenue par raisonnement, pas parce que les droites « ont l'air parallèles ».

\section{15. Enchaîner plusieurs propriétés}

Dans une configuration, on peut devoir utiliser plusieurs relations.

Supposons qu'un angle mesure :
\[
118^\circ.
\]

L'angle adjacent supplémentaire vaut :
\[
180^\circ-118^\circ=62^\circ.
\]

L'angle opposé par le sommet à cet angle mesure aussi :
\[
62^\circ.
\]

Le raisonnement comporte donc deux étapes distinctes.

\section{16. Démontrer plutôt que mesurer}

Au cycle 4, la mesure n'est plus une preuve.

Un angle mesuré à :
\[
60^\circ
\]
sur un dessin peut en réalité représenter une autre mesure si le schéma n'est pas à l'échelle.

\begin{attention}
Le rapporteur sert à construire ou à explorer. Une démonstration s'appuie sur les données et les propriétés.
\end{attention}

\section{17. Demi-tour et parallélisme}

Le demi-tour conserve le parallélisme.

Si une droite $(d)$ a pour image $(d')$ par un demi-tour et si le centre n'est pas sur $(d)$, alors :
\[
(d)\parallel(d').
\]

Cette propriété crée un lien entre transformation et géométrie des droites.

\section{18. Demi-tour et angles}

Comme le demi-tour conserve les mesures d'angles, l'image d'un angle de :
\[
35^\circ
\]
mesure :
\[
35^\circ.
\]

Cette invariance permet de raisonner sans reconstruire toutes les mesures.

\section{19. Exemple de construction}

On veut construire l'image d'un segment $[AB]$ par le demi-tour de centre $O$.

On construit d'abord :
\[
A'
\]
tel que $O$ soit le milieu de $[AA']$.

Puis :
\[
B'
\]
tel que $O$ soit le milieu de $[BB']$.

Le segment image est :
\[
[A'B'].
\]

On doit avoir :
\[
A'B'=AB.
\]

\section{20. Méthode pour déterminer un angle}

\begin{methode}
\begin{enumerate}
\item repérer les droites et la sécante ;
\item identifier la relation entre les angles ;
\item utiliser une propriété connue ;
\item écrire le calcul éventuel ;
\item conclure avec la mesure et l'unité degré.
\end{enumerate}
\end{methode}

\section{21. Méthode pour prouver le parallélisme}

\begin{methode}
\begin{enumerate}
\item identifier une paire d'angles alternes-internes ou correspondants ;
\item établir qu'ils ont la même mesure ;
\item citer la propriété caractéristique ;
\item conclure que les droites sont parallèles.
\end{enumerate}
\end{methode}

\section{22. Rédaction d'un raisonnement}

Une rédaction géométrique peut être concise.

Exemple :
\begin{quote}
Les angles considérés sont correspondants et ont tous deux pour mesure $64^\circ$. Donc les deux droites sont parallèles.
\end{quote}

Le point essentiel est d'indiquer :
\begin{itemize}
\item les faits connus ;
\item la propriété utilisée ;
\item la conclusion.
\end{itemize}

\section{23. Erreurs fréquentes}

\begin{itemize}
\item Confondre symétrie axiale et symétrie centrale.
\item Oublier que le centre est le milieu de $[AA']$.
\item Confondre angles alternes-internes et angles opposés par le sommet.
\item Déduire le parallélisme uniquement d'un dessin.
\item Utiliser une égalité d'angles sans identifier leur relation.
\item Mesurer au rapporteur alors qu'une propriété suffit.
\item Confondre une propriété directe et sa réciproque.
\end{itemize}

\section{À retenir}

\begin{aretenir}
Un demi-tour est une rotation de :
\[
\boxed{180^\circ}
\]
autour d'un centre.

Si $A'$ est l'image de $A$ de centre $O$ :
\[
\boxed{O\text{ est le milieu de }[AA']}.
\]

Le demi-tour conserve longueurs, angles et parallélisme.

Pour deux droites parallèles coupées par une sécante, les angles alternes-internes et correspondants associés sont égaux.

Réciproquement, l'égalité de ces angles permet de caractériser le parallélisme.
\end{aretenir}
